\documentclass[11pt,a4paper]{article}
\usepackage[T1]{fontenc}
\usepackage[margin=1in]{geometry}
\usepackage{amsmath,amssymb,amsthm}
\usepackage{hyperref}
\usepackage{booktabs}
\hypersetup{colorlinks=true,linkcolor=blue,citecolor=blue,urlcolor=blue}
\newtheorem{theorem}{Theorem}
\newtheorem{proposition}[theorem]{Proposition}
\newtheorem{corollary}[theorem]{Corollary}
\newtheorem{lemma}[theorem]{Lemma}
\newtheorem{definition}[theorem]{Definition}
\newtheorem{remark}{Remark}
\setlength{\emergencystretch}{1.5em}

\title{\textbf{Cost Self-Consistency and the Vacuum Einstein Equations:\\[2pt]
An Equivalence Theorem for Spherically Symmetric Spacetimes}}
\author{Y.\,Y.\,N.\ Li\footnote{Independent Researcher, China.
E-mail: \texttt{papaso@icloud.com}.
ORCID: \href{https://orcid.org/0009-0002-6471-139X}{0009-0002-6471-139X}.}}
\date{}

\begin{document}
\maketitle

\begin{abstract}
The $K{=}1$ Chronogeometrodynamics framework derives Lorentzian signature
from axioms on an information-time cost function $d(x;\delta x)$.
The point-level self-consistency condition $K = \boldsymbol{x}^T\!G\,\boldsymbol{x} = 1$
encodes the balance between temporal and spatial cost.
We introduce two field-level analogues for static spherically symmetric
spacetimes $ds^2 = -f\,dt^2 + f^{-1}dr^2 + r^2 d\Omega^2$:
the radial cost $K_{\mathrm{f}} \equiv \sigma_2^2\,\Box\ln\sigma_1
= f + 2rf' + \frac{r^2}{2}f''$
and the angular cost $K_{\mathrm{a}} \equiv rf'+f$.
Our main result is an algebraic equivalence:
\[
  K_{\mathrm{f}} = 1 \;\wedge\; K_{\mathrm{a}} = 1
  \quad\Longleftrightarrow\quad R_{\mu\nu} = 0\,.
\]
The radial condition $K_{\mathrm{f}}=1$ alone gives the Ricci scalar equation
$R=0$, admitting the two-parameter family $f = 1 - C_1/r - C_2/r^2$;
the angular condition $K_{\mathrm{a}}=1$ selects Schwarzschild uniquely.
In the weak field, $K_{\mathrm{a}}=1$ yields the Newtonian potential
$\varphi = C/r$ directly.
While the radial observable has a clear cost origin, the angular condition
is defined to match $R_{\theta\theta}=0$ rather than derived from cost axioms.
The extension to general metrics parallels Jacobson's null-direction patching.
We discuss these limitations, the cost interpretation of matter, and the
prospect that a minimum symplectic eigenvalue $\sigma_{1,\min}>0$
could bound curvature and resolve singularities.
\end{abstract}

%===================================================================
\section{Introduction}
\label{sec:intro}
%===================================================================

\subsection{Context}

Several derivations of Einstein's field equations exist beyond Einstein's
original geometric argument~\cite{Einstein1915} and
Hilbert's variational principle~\cite{Hilbert1915}.
Jacobson~\cite{Jacobson1995} showed that applying the Clausius
relation $\delta Q = T\,\delta S$ with the Bekenstein--Hawking
entropy $S = A/4$ on every local Rindler horizon yields
$G_{\mu\nu} = 8\pi T_{\mu\nu}$.
Padmanabhan~\cite{Padmanabhan2010} and Verlinde~\cite{Verlinde2011}
proposed related thermodynamic and entropic approaches.
All assume Lorentzian signature as a starting point.

The $K{=}1$ Chronogeometrodynamics framework~\cite{Li2026math,Li2026real}
removes this assumption.
Three axioms on an information-time cost function---Realizability~(R),
Expansion~(E), and Temporal cost~(T)---fix the Hessian signature
to $(1,1)$, deriving the Lorentzian character of the metric.
The self-consistency condition $K = \boldsymbol{x}^T\!G\,\boldsymbol{x} = 1$
drives Ornstein--Uhlenbeck dynamics whose stationary variance
reproduces the Tolman--Unruh temperature, furnishing two of
Jacobson's three inputs.
The Einstein equations then follow via the Jacobson bridge once
$S \propto A$ is imposed.

\subsection{Motivation and scope}

A natural question is whether the thermodynamic bridge is necessary,
or whether cost self-consistency can constrain the curvature more directly.
In this paper we investigate this question for the spherically symmetric sector.

We introduce two field-level cost observables---$K_{\mathrm{f}}$
and $K_{\mathrm{a}}$---and prove that their simultaneous condition
$K_{\mathrm{f}} = K_{\mathrm{a}} = 1$ is algebraically equivalent
to $R_{\mu\nu} = 0$.
However, three important caveats apply:
\begin{enumerate}
\item[(i)] The angular observable $K_{\mathrm{a}}$ is \emph{defined}
  to coincide with the GR condition $R_{\theta\theta}=0$;
  we do not derive it from the cost axioms R, E, T.
\item[(ii)] The extension to general metrics requires applying
  $K_{\mathrm{f}}=1$ on each null direction and assembling the
  tensor equation via the polarization identity---the same logical
  structure as Jacobson's argument with $K{=}1$ replacing Clausius.
\item[(iii)] In the vacuum sector the equivalence is exact,
  so \emph{no} observable deviation from GR is predicted.
\end{enumerate}
The contribution is therefore a cost-based \emph{reinterpretation}
of the vacuum Einstein equations, equipped with new observables and
a precise point-to-field correspondence, rather than an independent
derivation.

%===================================================================
\section{Review: point-level $K{=}1$}
\label{sec:point}
%===================================================================

We recall the essentials from~\cite{Li2026math}.
On a local Rindler wedge the $2\times 2$ cost metric is
$G = \mathrm{diag}(-\sigma_1^2,\,1)$, where $\sigma_1$ encodes the
acceleration scale.
The self-consistency condition
\begin{equation}\label{eq:K1}
  K \;=\; \sigma_1^2\,x_0^2 - x_1^2 \;=\; 1
\end{equation}
defines a hyperboloid: temporal cost equals spatial cost.
A smooth restoring potential $V = \tfrac{1}{2}(K{-}1)^2$
generates OU dynamics with stationary variance
$\mathrm{Var}(K) = T_{\mathrm{eff}}$, identified with the
Tolman temperature.

At the field level, one expects an analogous balance condition on the
metric itself.
The metric is parameterized by two symplectic eigenvalues
$\sigma_1(r)$ and $\sigma_2(r)$, and the field equations should
constrain how these eigenvalues vary across spacetime.

%===================================================================
\section{Field-level observables}
\label{sec:field}
%===================================================================

We work with the static spherically symmetric metric
\begin{equation}\label{eq:metric}
  ds^2 = -f(r)\,dt^2 + f(r)^{-1}\,dr^2 + r^2\,d\Omega^2\,,
\end{equation}
with $f(r) > 0$ outside the horizon.
The symplectic eigenvalues are
\begin{equation}\label{eq:sigmas}
  \sigma_1 = r\sqrt{f}\,,\qquad \sigma_2 = r\,.
\end{equation}

\begin{definition}[Radial cost]\label{def:Kf}
The radial--temporal cost observable is
\begin{equation}\label{eq:Kf}
  K_{\mathrm{f}} \;:=\; \sigma_2^2\;\Box\ln\sigma_1\,,
\end{equation}
where $\Box\phi = (r^2 f)^{-1}\,\partial_r(r^2 f\,\partial_r\phi)$
is the scalar d'Alembertian of~\eqref{eq:metric}.
\end{definition}

Writing $\ln\sigma_1 = \ln r + \frac{1}{2}\ln f$ and computing:
\begin{equation}\label{eq:Kf-explicit}
  K_{\mathrm{f}} \;=\; f + 2r\,f' + \tfrac{r^2}{2}\,f''\,.
\end{equation}

The physical content of $K_{\mathrm{f}}$ is transparent.
The identity $R = -2\,\Box\ln\sigma_1 + 2/r^2$
(established in~\cite{Li2026symp}) gives
\begin{equation}\label{eq:KfR}
  K_{\mathrm{f}} = 1 \quad\Longleftrightarrow\quad
  \Box\ln\sigma_1 = \frac{1}{r^2} = \frac{1}{\sigma_2^2}
  \quad\Longleftrightarrow\quad R = 0\,.
\end{equation}
The radial variation of $\sigma_1$ balances the angular curvature
cost $1/\sigma_2^2$ precisely when the Ricci scalar vanishes.
This extends the point-level balance (temporal = spatial) to the
field level (radial variation cost = angular curvature cost).

\begin{definition}[Angular cost]\label{def:Ka}
The angular cost observable is
\begin{equation}\label{eq:Ka}
  K_{\mathrm{a}} \;:=\; r\,f' + f \;=\; \frac{d(rf)}{dr}\,.
\end{equation}
\end{definition}

\begin{remark}[Status of $K_{\mathrm{a}}$]\label{rem:Ka}
The standard GR result $R_{\theta\theta} = 1 - f - rf'$ shows that
$K_{\mathrm{a}}=1$ is simply $R_{\theta\theta}=0$.
Equivalently, $K_{\mathrm{a}}=1$ states that the Misner--Sharp
mass $m(r) = \frac{r}{2}(1-f)$ is constant: $dm/dr=0$ (no local
mass source).

Unlike $K_{\mathrm{f}}$, the angular observable $K_{\mathrm{a}}$ is
\emph{not derived} from the cost axioms;
it is \emph{defined} to match $R_{\theta\theta}=0$.
Deriving $K_{\mathrm{a}}=1$ from the information-time cost structure
remains an open problem (see Section~\ref{sec:open}).
This is the primary conceptual gap in the present construction.
\end{remark}

%===================================================================
\section{Equivalence theorem}
\label{sec:theorem}
%===================================================================

\begin{theorem}\label{thm:main}
For the metric~\eqref{eq:metric} with $f\in C^2(0,\infty)$:
\begin{equation}\label{eq:main}
  K_{\mathrm{f}} = 1 \;\wedge\; K_{\mathrm{a}} = 1
  \quad\Longleftrightarrow\quad
  R_{\mu\nu} = 0\,.
\end{equation}
\end{theorem}

\begin{proof}
The independent Ricci components for~\eqref{eq:metric} are
\begin{equation}\label{eq:Ricci}
  R_{tt} \propto f'' + \frac{2f'}{r}\,,\qquad
  R_{\theta\theta} = 1 - f - rf'\,.
\end{equation}
($R_{rr}=0$ is equivalent to $R_{tt}=0$ for this diagonal ansatz.)

\medskip\noindent
$(\Rightarrow)$\;
$K_{\mathrm{a}}=1$ gives $rf'+f=1$, i.e.\ $R_{\theta\theta}=0$.
Subtracting: $K_{\mathrm{f}} - K_{\mathrm{a}} = rf' + \frac{r^2}{2}f'' = 0$,
hence $f'' = -2f'/r$ and $R_{tt}=0$.

\medskip\noindent
$(\Leftarrow)$\;
$R_{\theta\theta}=0$ gives $K_{\mathrm{a}}=1$.
$R_{tt}=0$ gives $f''=-2f'/r$.
Substituting:
$K_{\mathrm{f}} = f + 2rf' + \frac{r^2}{2}(-\frac{2f'}{r})
= f + rf' = 1$,
using $K_{\mathrm{a}}=1$.
\end{proof}

\begin{corollary}\label{cor:scalar}
$K_{\mathrm{f}}=1$ is equivalent to $R=0$ (Ricci \emph{scalar}),
which is strictly weaker than $R_{\mu\nu}=0$ (Ricci \emph{tensor}).
The angular condition $K_{\mathrm{a}}=1$ provides the additional
content needed to promote the scalar equation to a full tensor equation.
\end{corollary}

%===================================================================
\section{Solution structure}
\label{sec:solutions}
%===================================================================

\subsection{General solution of $K_{\mathrm{f}}=1$}

The ODE
\begin{equation}\label{eq:ode}
  f + 2r\,f' + \tfrac{r^2}{2}\,f'' = 1
\end{equation}
admits a power-law ansatz $f = 1 - A\,r^{-n}$.
The coefficient of $A\,r^{-n}$ is $-(n{-}1)(n{-}2)/2$,
which vanishes for $n=1$ and $n=2$.
The general solution is
\begin{equation}\label{eq:gensol}
  f(r) = 1 - \frac{C_1}{r} - \frac{C_2}{r^2}\,,
\end{equation}
a two-parameter family.
With $C_1 = 2M$ and $C_2 = Q^2$ this resembles the
Reissner--Nordstr\"om metric, but here both parameters arise as
integration constants of the \emph{scalar} equation $R=0$.

\subsection{Role of $K_{\mathrm{a}}=1$}

For~\eqref{eq:gensol}, $K_{\mathrm{a}} = 1 + C_2/r^2$.
The angular condition forces $C_2 = 0$, and with $C_1 = 2M$
one recovers Schwarzschild uniquely---Birkhoff's theorem.

\subsection{Linearized regime}

For $f = 1 + \varepsilon\,\varphi(r)$ with $|\varepsilon|\ll 1$:
\begin{equation}\label{eq:lin}
  K_{\mathrm{a}} = 1:\quad
  r\,\varphi' + \varphi = 0
  \;\;\Longrightarrow\;\;
  \frac{d(r\varphi)}{dr} = 0
  \;\;\Longrightarrow\;\;
  \varphi = \frac{C}{r}\,.
\end{equation}
With $C = -2M$ this is the Newtonian potential.
The radial condition $K_{\mathrm{f}}=1$ is automatically
satisfied at linear order: $\varphi + 2r\varphi' + \frac{r^2}{2}\varphi''
= C/r - 2C/r + C/r = 0$.

This means $K_{\mathrm{f}}$ carries independent information only at
\emph{nonlinear} order, i.e., in the strong-field regime.
In the weak field, the angular condition alone determines the solution.

%===================================================================
\section{Non-vacuum sector}
\label{sec:nonvac}
%===================================================================

\subsection{Matter as cost imbalance}

The Einstein $G_{tt}$ component reads
$(1{-}f)/r^2 - f'/r = 8\pi\rho$,
which rearranges to
\begin{equation}\label{eq:rho}
  \rho \;=\; \frac{1 - K_{\mathrm{a}}}{8\pi\,r^2}\,.
\end{equation}
The energy density equals the angular cost deficit per unit proper area.
Normal matter ($\rho > 0$) corresponds to $K_{\mathrm{a}} < 1$:
the metric fails to achieve angular self-consistency, and the failure
is measured by $1 - K_{\mathrm{a}}$.

This has been verified symbolically for both constant density
$\rho = \rho_0$ (interior Schwarzschild) and variable density
$\rho = \rho_0(1-r/a)$ profiles.

We stress that~\eqref{eq:rho} is a \emph{rewriting} of the
Einstein equations, not a derivation from cost alone.
The coefficient $8\pi$ is inherited from GR; the cost framework
does not predict it independently.

\subsection{Cosmological constant}

For the Schwarzschild--de~Sitter metric
$f = 1 - 2M/r - \Lambda r^2/3$:
\begin{equation}\label{eq:Lambda}
  K_{\mathrm{a}} = 1 - \Lambda\,r^2\,,\qquad
  K_{\mathrm{f}} = 1 - 2\Lambda\,r^2\,.
\end{equation}
The cosmological constant shifts the self-consistency target;
notably, $K_{\mathrm{f}}$ and $K_{\mathrm{a}}$ are shifted by
different amounts ($2\Lambda r^2$ vs $\Lambda r^2$).
The value of $\Lambda$ is not predicted by the framework.

%===================================================================
\section{General metrics and the Jacobson parallel}
\label{sec:general}
%===================================================================

\subsection{Polarization identity}

The extension beyond spherical symmetry relies on the following result.

\begin{proposition}\label{prop:polar}
Let $R_{\mu\nu}$ be a symmetric tensor on an $n$-dimensional
Lorentzian manifold with $R = g^{\mu\nu}R_{\mu\nu} = 0$.
If $R_{\mu\nu}\,v^\mu v^\nu = 0$ for every null $v^\mu$,
then $R_{\mu\nu} = 0$.
\end{proposition}

\begin{proof}
Choose an orthonormal frame $\{e_0,e_1,\ldots,e_{n-1}\}$ with
$e_0$ timelike.
Null vectors are $v = e_0 + n^i e_i$ with $|n|=1$.
Taking $n = \pm e_k$:
$R_{00} + R_{kk} \pm 2R_{0k} = 0$,
giving $R_{0k}=0$ and $R_{00}+R_{kk}=0$ for each $k$.
Taking $n = (e_i{+}e_j)/\sqrt{2}$:
$R_{00} + (R_{ii}{+}2R_{ij}{+}R_{jj})/2 = 0$;
using $R_{ii}=R_{jj}=-R_{00}$ gives $R_{ij}=0$ for $i\neq j$.
Thus $R_{\mu\nu} = R_{00}\,\mathrm{diag}(-1,\ldots,-1)$.
The trace $R = -nR_{00} = 0$ forces $R_{00}=0$.
\end{proof}

\subsection{Proof sketch for general metrics}

For a general metric, at each point $p$ and each null direction $v^\mu$
one can define a local Rindler wedge with its own symplectic eigenvalue
$\sigma_1(p,v)$.
The cost condition $K_{\mathrm{f}}(p,v) = 1$ is then asserted
for every null direction.
If this implies $R_{\mu\nu}v^\mu v^\nu = 0$
(as it does in the spherically symmetric case via the explicit
formula~\eqref{eq:KfR}), then combined with $R=0$ the
polarization identity gives $R_{\mu\nu}=0$.

\subsection{The missing step and the Jacobson parallel}

The implication $K_{\mathrm{f}}(v)=1 \Rightarrow R_{\mu\nu}v^\mu v^\nu=0$
for general metrics requires the Gauss--Codazzi relations for null
hypersurfaces connecting the intrinsic 2D curvature to the
4D Ricci focusing $R_{\mu\nu}v^\mu v^\nu$.
This step is \emph{not yet proved}.

We note that the logical architecture of this general proof is
identical to Jacobson's~\cite{Jacobson1995}:
a local condition (here $K_{\mathrm{f}}=1$; in Jacobson's case, the
Clausius relation) is applied on each null direction, and the
full tensor equation is assembled by polarization.
The \emph{content} differs---information-time cost replaces
thermodynamic heat---but the \emph{patching logic} is shared.
This parallel limits the extent to which the present approach
can claim independence from Jacobson's construction.

%===================================================================
\section{Point--field correspondence}
\label{sec:correspondence}
%===================================================================

The structural parallel is summarized in Table~\ref{tab:corr}.

\begin{table}[h]
\centering
\begin{tabular}{@{}lll@{}}
\toprule
& \textbf{Point level (2D)} & \textbf{Field level (4D)} \\
\midrule
Observable & $K = \boldsymbol{x}^{T}\!G\,\boldsymbol{x}$
  & $K_{\mathrm{f}} = \sigma_2^2\,\Box\ln\sigma_1$;\;
    $K_{\mathrm{a}} = rf'+f$ \\[3pt]
Balance & temporal = spatial & radial variation = angular curvature \\[3pt]
Self-consistency & $K = 1$ & $K_{\mathrm{f}}{=}1\;\wedge\;K_{\mathrm{a}}{=}1
  \;\Leftrightarrow\;R_{\mu\nu}{=}0$ \\[3pt]
Departure & $K \neq 1 \to V > 0$ & $K_{\mathrm{a}}\neq 1\;\to\;\rho = (1{-}K_{\mathrm{a}})/(8\pi r^2)$ \\[3pt]
Potential & $V = \tfrac{1}{2}(K{-}1)^2$ &
  $V_{\mathrm{f}} = \tfrac{1}{2}(K_{\mathrm{f}}{-}1)^2$ \\[3pt]
Conditions & 1 (2D wedge) & 2 (4D: radial + angular) \\
\bottomrule
\end{tabular}
\caption{Point-level $K{=}1$ versus field-level cost self-consistency.
At the point level one condition suffices because the Rindler wedge is
two-dimensional; in 4D each independent geometric sector contributes
one condition.}
\label{tab:corr}
\end{table}

%===================================================================
\section{Toward singularity resolution}
\label{sec:singularity}
%===================================================================

The most physically consequential open question in the $K{=}1$
program may be whether the symplectic eigenvalue $\sigma_1$ admits
a positive lower bound.
If $K$ and the boost angle $\theta$ are canonically conjugate
(as conjectured in Open Question~7 of~\cite{Li2026math}),
the uncertainty relation $\sigma_1\,\Delta\theta \geq \hbar/2$
would imply $\sigma_1 \geq \sigma_{1,\min} > 0$.

A minimum $\sigma_1$ yields a maximum curvature.
In the Schwarzschild interior, $R_{\mu\nu\rho\sigma}R^{\mu\nu\rho\sigma}
= 48M^2/r^6$ diverges as $r\to 0$.
If $\sigma_1 \geq \sigma_{1,\min} \sim \ell_P/2$, then:
\begin{equation}
  A_{\min} = 4\pi\sigma_{1,\min}^2 \sim \pi\,\ell_P^2\,,\qquad
  R_{\max} \sim \frac{2}{\sigma_{1,\min}^2} \sim 3\times 10^{70}\;\mathrm{m^{-2}}\,,
\end{equation}
bounding curvature at the Planck scale and replacing the singularity
with a Planck-density core.
For comparison, the curvature at the horizon of a solar-mass
black hole is $\sim 2\times 10^{-7}\;\mathrm{m^{-2}}$, a factor of
$\sim 10^{77}$ below $R_{\max}$.

Proving $\sigma_{1,\min}>0$ would constitute a singularity-resolution
result within the $K{=}1$ framework, without invoking quantum gravity.
This remains the most promising direction for future work.

%===================================================================
\section{Discussion}
\label{sec:discussion}
%===================================================================

\subsection{Summary of results}

\begin{enumerate}
\item We defined field-level cost observables $K_{\mathrm{f}}$
  and $K_{\mathrm{a}}$ for the spherically symmetric sector.
\item We proved the exact equivalence
  $K_{\mathrm{f}}{=}1 \wedge K_{\mathrm{a}}{=}1
  \Leftrightarrow R_{\mu\nu}{=}0$
  (Theorem~\ref{thm:main}), including both directions.
\item We derived the general solution of $K_{\mathrm{f}}{=}1$
  ($f = 1 - C_1/r - C_2/r^2$) and showed that $K_{\mathrm{a}}{=}1$
  selects Schwarzschild.
\item We obtained the Newtonian potential from $K_{\mathrm{a}}{=}1$
  in one line (Eq.~\ref{eq:lin}), and showed that $K_{\mathrm{f}}$
  is redundant at linear order.
\item We interpreted matter as cost imbalance~\eqref{eq:rho}
  and identified the cosmological constant as a target shift~\eqref{eq:Lambda}.
\end{enumerate}

\subsection{Limitations}

\begin{enumerate}
\item \textbf{$K_{\mathrm{a}}$ is defined, not derived.}
  The angular condition $K_{\mathrm{a}}=1$ is imposed to match
  $R_{\theta\theta}=0$; we lack a cost derivation.
  This is the primary conceptual gap.
\item \textbf{Spherical symmetry only.}
  The equivalence theorem is proved for the metric ansatz~\eqref{eq:metric}.
  The general proof requires the null-surface Gauss--Codazzi step,
  which is not completed.
\item \textbf{Jacobson parallel.}
  The general proof has the same logical structure as Jacobson's
  thermodynamic derivation: a local condition on each null direction,
  assembled by polarization.
  The content differs ($K{=}1$ vs Clausius), but the architecture is shared.
\item \textbf{No new predictions.}
  In the vacuum sector the equivalence is exact, so no deviation from
  GR is predicted.
  The matter interpretation uses Einstein's equations as input.
  These limitations are shared by all known ``derivations'' of Einstein's
  equations~\cite{Jacobson1995,Padmanabhan2010,Verlinde2011}.
\item \textbf{Linearized redundancy.}
  $K_{\mathrm{f}}$ is automatically satisfied at linear order;
  it carries independent content only in the strong-field regime.
  This limits any future testable consequences to strong-gravity settings.
\end{enumerate}

\subsection{What is genuinely new}

Despite the limitations above, the following elements are new:
\begin{itemize}
\item The \emph{observables} $K_{\mathrm{f}}$ and $K_{\mathrm{a}}$
  as field-level cost quantities with concrete formulas.
\item The \emph{equivalence theorem} (Theorem~\ref{thm:main}):
  a precise, verifiable statement with a complete two-directional proof.
\item The \emph{ODE}~\eqref{eq:ode} and its closed-form general solution,
  exhibiting the scalar-vs-tensor distinction ($R{=}0$ vs $R_{\mu\nu}{=}0$).
\item The \emph{cost interpretation} of energy density:
  $\rho = (1{-}K_{\mathrm{a}})/(8\pi r^2)$,
  reading matter as the angular cost deficit.
\item The \emph{point--field correspondence} (Table~\ref{tab:corr}),
  extending the $K{=}1$ framework from horizons to spacetime.
\end{itemize}

\subsection{Open problems}
\label{sec:open}

\begin{enumerate}
\item Derive $K_{\mathrm{a}}=1$ from the cost axioms R, E, T.
\item Complete the general proof (null-surface Gauss--Codazzi).
\item Prove $\sigma_{1,\min}>0$ and establish curvature bounds.
\item Identify testable predictions in the strong-field regime.
\end{enumerate}

\section*{Acknowledgments}

All symbolic results were verified with \texttt{sympy}; numerical
checks with \texttt{numpy}.
A supplementary test suite is available containing 51 formula verifications
and 79 reasoning checks with 28 explicitly flagged caveats.

\begin{thebibliography}{99}

\bibitem{Einstein1915}
A.~Einstein,
``Die Feldgleichungen der Gravitation,''
\textit{Sitzungsber.\ Preuss.\ Akad.\ Wiss.\ Berlin}, 844 (1915).

\bibitem{Hilbert1915}
D.~Hilbert,
``Die Grundlagen der Physik,''
\textit{Nachr.\ Ges.\ Wiss.\ G\"ottingen}, 395 (1915).

\bibitem{Jacobson1995}
T.~Jacobson,
``Thermodynamics of spacetime: The Einstein equation of state,''
\textit{Phys.\ Rev.\ Lett.}\ \textbf{75}, 1260 (1995).

\bibitem{Padmanabhan2010}
T.~Padmanabhan,
``Thermodynamical aspects of gravity: New insights,''
\textit{Rep.\ Prog.\ Phys.}\ \textbf{73}, 046901 (2010).

\bibitem{Verlinde2011}
E.~Verlinde,
``On the origin of gravity and the laws of Newton,''
\textit{J.\ High Energy Phys.}\ \textbf{2011}(04), 029 (2011).

\bibitem{Li2026math}
Y.\,Y.\,N.~Li,
``$K{=}1$ Chronogeometrodynamics: Lorentzian geometry from information time,''
submitted to \textit{Found.\ Phys.} (2026).

\bibitem{Li2026real}
Y.\,Y.\,N.~Li,
``Realizability and the origin of causality,''
submitted to \textit{Found.\ Phys.} (2026).

\bibitem{Li2026symp}
Y.\,Y.\,N.~Li,
``Symplectic gravity: from Rindler horizons to Einstein equations,''
in preparation (2026).

\end{thebibliography}

\end{document}
